\documentclass[journal]{IEEEtran}
\usepackage{cite}
\usepackage{graphicx}
\usepackage{amsmath,amssymb,amsfonts}
\usepackage{algorithmic}
\usepackage{url}
\usepackage{hyperref}

\begin{document}

\title{OrganoidEnv: A Stabilized Reinforcement Learning Environment for Training Biologically Constrained Spiking Neural Networks via Sensory--Motor Mapping}

\author{Vansh~Sharma%
\thanks{V. Sharma is with the Department of Electronics and Communication Engineering, Vivekananda Institute Of Professional Studies, New Delhi, India (e-mail: vansh.sharma@vips.edu).}
}

\maketitle

\begin{abstract}
Integrating biologically realistic spiking neural networks (SNNs) with modern reinforcement learning (RL) remains challenging because of the non-differentiable nature of spikes, recurrent instability of spiking dynamics, and metabolic constraints that are typically ignored in artificial neural networks (ANNs).
This study introduces \emph{OrganoidEnv}, a Gymnasium-compatible environment that wraps a 500-neuron metabolic-Izhikevich network and exposes it as a controllable in silico ``organoid'' to a deep RL agent.
Unlike approaches that directly modify SNN synaptic weights via surrogate gradients, OrganoidEnv employs an outer Duelling Double Deep Q-Network (D3QN) that interacts with the SNN solely through current-based stimulation of the sensory and motor pathways.
Learning within the SNN is governed by a dopamine D3-inspired dual-trace eligibility rule, homeostatic inhibitory scaling, structural plasticity, and a Global Activity Regulator (GAR) that maintains activity within a biologically realistic firing-rate regime.
Using a four-stage curriculum in a continuous 2D navigation task with obstacles and context-dependent multi-goal switching, OrganoidEnv attained an overall success rate of 70.8\% across 500 episodes, including 97.0\% success in obstacle navigation.
Comparisons with ANN baselines demonstrate that ANNs optimized purely for performance naturally outperform biologically constrained systems; however, OrganoidEnv's goal is not performance parity but biological fidelity and interpretability as a neuromorphic testbed.
\end{abstract}

\begin{IEEEkeywords}
Spiking Neural Networks, Reinforcement Learning, Neuromorphic Computing, Organoid Intelligence, Sensory--Motor Mapping, Synaptic Plasticity
\end{IEEEkeywords}

\section{Introduction}

\IEEEPARstart{S}{piking} neural networks (SNNs) are a central abstraction for neuromorphic computation due to their event-driven communication, temporal coding, and close correspondence to biological neurons and synapses \cite{Izhikevich2003simple,NeuromorphicReview2019}. 
In contrast, artificial neural networks (ANNs) trained with backpropagation deliver state-of-the-art performance across domains but typically ignore biophysical constraints, such as sparse firing, local plasticity, and metabolic limits \cite{Mnih2015human}. 
Bridging the gap between high-performing deep RL and biologically grounded SNNs remains a challenge.

Classical spike-timing-dependent plasticity (STDP) and reward-modulated STDP (R-STDP) provide local learning rules but struggle with temporal credit assignment in delayed-reward tasks, especially in recurrent networks \cite{Florian2007rstdp,FarriesFairhall2007}. 
Surrogate-gradient methods and e-prop approximate backpropagation through time for spiking systems \cite{Bellec2018l2l,DirectTrainingSNN2024}, yet they often sacrifice biological plausibility by relying on continuous relaxations and storing long histories of membrane state.

In parallel, recent work on \emph{organoid intelligence} (OI) has proposed the combination of living brain organoids with machine interfaces to create hybrid biocomputing systems \cite{Smirnova2023OI,Hartung2023OIworkshop}. 
Although in vitro organoids provide biological realism, they are expensive, fragile, and experimentally constrained. 
There is a need for in silico organoid analogues that capture the key principles of organoid learning and plasticity while being reproducible and controllable.

This study proposes \emph{OrganoidEnv}, a stabilised RL environment that treats a biologically constrained SNN as an \emph{actor embedded in an environment}, rather than as a differentiable function to be optimised end-to-end. 
An external RL agent plays the role of an automated experimentalist that learns how to stimulate the SNN to elicit the desired motor behaviours, while plasticity within the SNN remains local and dopamine-modulated.

\subsection{Contributions}

The main contributions of this study are as follows:

\begin{itemize}
    \item \textbf{Organoid-like SNN wrapper:} A Gymnasium-compatible environment built around a 500-neuron Metabolic-Izhikevich network with ATP-like energy variables, capturing metabolic exhaustion and recovery.
    \item \textbf{Biological stabilizers:} A global activity regulator (GAR), inhibitory homeostasis, and activity-dependent structural plasticity (synapse pruning) that jointly prevent seizure and coma states.
    \item \textbf{Dopamine dual-trace rule:} A D3-inspired dual eligibility-trace mechanism that combines millisecond-scale STDP with second-scale reward-modulated traces to address delayed credit assignment without surrogate gradients.
    \item \textbf{Sensory--motor mapping:} A direct motor-quadrant stimulation scheme that enables a deep RL agent to train the SNN via current injections while the SNN itself maintains local plasticity and homeostasis.
    \item \textbf{Experimental evaluation:} A six-stage curriculum on a continuous 2D navigation task with obstacles and context-dependent goals, including ablations and ANN baselines to quantify the role of stabilizers and dual-trace plasticity.
\end{itemize}

\section{Related Work}

\subsection{Spiking neural networks and neuromorphic learning}

SNNs have been widely studied as candidates for energy-efficient computation and brain-inspired learning \cite{Izhikevich2003simple,NeuromorphicSNNMemristor2019,NeuromorphicReview2019,Davies2024Loihi}. 
Direct training of deep SNNs using surrogate gradients or spike-based approximations has recently achieved competitive performance on standard benchmarks \cite{RecentFrontiersSNN2022,DirectTrainingSNN2024}. 
However, these approaches often relax spiking nonlinearity and rely on global gradient information that is difficult to reconcile with strict locality and metabolic constraints.

Homeostatic plasticity and synaptic scaling provide mechanisms to stabilise firing rates and maintain excitation--inhibition (E/I) balance in recurrent networks \cite{Turrigiano2008SelfTuning,HomeostaticSTDP2010,eLifeHomeostatic2025}. 
Structural plasticity, including synaptogenesis and pruning, is known to shape network topology over longer timescales but is less frequently incorporated into algorithmic SNN frameworks.

\subsection{Reward-modulated STDP and RL--SNN hybrids}

Reward-modulated STDP extends classical STDP by combining local spike timing with global reward signals, enabling simple reinforcement learning in spiking networks \cite{Florian2007rstdp,FarriesFairhall2007}. 
Nevertheless, scaling R-STDP to complex temporally extended tasks remains challenging because of sparse reward signals and long credit assignment horizons.

Hybrid RL--SNN architectures often use ANN controllers to directly adjust SNN synapses or provide modulatory signals to SNNs \cite{AcetylcholineNav2018,Bellec2018l2l,Zheng2024SNNRL}. 
Many of these methods aim for performance rather than biological plausibility and rarely incorporate explicit metabolic or homeostatic constraints.

\subsection{Organoid intelligence and in silico organoids}

Organoid intelligence has emerged as a new frontier in biocomputing, leveraging living brain organoids interfaced with electrodes and AI systems \cite{Smirnova2023OI,Hartung2023OIworkshop,Xiong2024BridgingAI}. 
These studies argue for a hybrid paradigm in which AI algorithms guide the experimental stimulation of organoids, whereas the organoids themselves adapt via biological plasticity. 
OrganoidEnv can be viewed as an in silico counterpart of this paradigm, where an RL agent plays the role of an automated experimentalist, controlling a simulated organoid-like SNN.

\section{Architecture and Methods}

OrganoidEnv consists of three coupled components: (1) a 500-neuron metabolic-Izhikevich SNN, (2) biological stabilisers (GAR, homeostasis, and structural plasticity), and (3) an outer Duelling Double DQN agent that interacts with the SNN via current injections and reads out motor activity.

\subsection{Metabolic-Izhikevich neuron model}

Each neuron $i$ follows the Izhikevich model augmented with a scalar energy variable $E_i$ which represents ATP availability. 
The subthreshold dynamics are as follows:
\begin{align}
\frac{dv_i}{dt} &= 0.04 v_i^2 + 5 v_i + 140 - u_i + I_i, \label{eq:izh_v} \\
\frac{du_i}{dt} &= a (b v_i - u_i),
\end{align}
where $v_i$ is the membrane potential, $u_i$ is the recovery variable, and $I_i$ is the total synaptic plus injected current. 
When $v_i \geq 30$~mV, the neuron emits a spike and is reset via $v_i \leftarrow c$ and $u_i \leftarrow u_i + d$ \cite{Izhikevich2003simple}.

The metabolic extension updates energy as follows:
\begin{equation}
E_i(t+1) = \min\left(E_{\max},\, E_i(t) + \frac{\Delta t}{\tau_{\mathrm{recovery}}}\right) - \delta_{\mathrm{spike}},
\label{eq:energy}
\end{equation}
where $E_{\max}$ is the maximum energy capacity, $\tau_{\mathrm{recovery}}$ is the ATP recovery time constant, and $\delta_{\mathrm{spike}}$ is the energy cost per spike. 
Spikes are suppressed whenever $E_i \leq 0$, modelling the metabolic fatigue.

To prevent numerical divergence under large injected currents, the membrane potentials were clamped after each integration step:
\begin{equation}
v_i^{\mathrm{clipped}} = \mathrm{clip}(v_i, v_{\min}, v_{\max}),
\end{equation}
with $v_{\min} = -200$~mV and $v_{\max} = 100$~mV. 
This preserves the canonical Izhikevich spiking regimes while avoiding NaNs under strong stimulation.

\subsection{Network topology}

The SNN comprises three functional layers.
\begin{itemize}
    \item \textbf{Sparse Distributed Memory (SDM) layer (256 neurons):} Receives a low-dimensional observation vector (proximity sensors, goal coordinates, context bit) and projects it into a 10\%-sparse high-dimensional representation, inspired by cerebellar granule cells and SDM principles.
    \item \textbf{Hidden layer (144 neurons):} A recurrent integration hub maintaining temporal context and supporting attractor-like dynamics.
    \item \textbf{Motor layer (100 neurons):} Split into four contiguous 25-neuron quadrants corresponding to Up, Down, Left, and Right motor directions.
\end{itemize}

Connectivity is dense and recurrent within layers and sparse between layers, with separate excitatory and inhibitory neuronal populations. 
Synaptic weights are initialised from zero-mean Gaussian distributions and then shaped by the plasticity and stabilisers. 
The key SNN-related hyperparameters are listed in Table~\ref{table:hyperparameters}.

\begin{figure}
    \centering
    \includegraphics[width=0.75\linewidth]{Metabolic-Izhikevich SNN-2026-03-27-134519.png}
    \caption{System architecture of OrganoidEnv. The outer Duelling Double DQN agent (top) interacts with the biologically constrained SNN (bottom) via current-based stimulation of sensory and motor pathways. The Global Activity Regulator, homeostatic inhibition, and structural plasticity jointly maintain network stability.}
    \label{fig:placeholder}
\end{figure}

\subsection{Dopamine D3 dual-trace learning rule}

To address the delayed reward credit assignment without backpropagation, each synapse $w_{ij}$ maintains a fast and slow eligibility trace:
\begin{align}
\frac{de_{ij}^{\mathrm{fast}}}{dt} &= -\frac{e_{ij}^{\mathrm{fast}}}{\tau_{\mathrm{fast}}} + \mathrm{STDP}(t_i^{\mathrm{pre}}, t_j^{\mathrm{post}}), \label{eq:fast} \\
\frac{de_{ij}^{\mathrm{slow}}}{dt} &= -\frac{e_{ij}^{\mathrm{slow}}}{\tau_{\mathrm{slow}}} + e_{ij}^{\mathrm{fast}}, \label{eq:slow}
\end{align}
with $\tau_{\mathrm{fast}} = 100$ms and $\tau_{\mathrm{slow}} = 2500$ms. 
The STDP term implements an exponential pre--post window with separate time constants for causal and anti-causal spike pairs, producing positive eligibility for pre-before-post events and negative eligibility for post-before-pre \cite{Florian2007rstdp}.

Whenever the environment emits a scalar reward $r$ (for goal attainment or penalty), the synaptic weights are updated as follows:
\begin{equation}
\Delta w_{ij} = \eta\, r\, e_{ij}^{\mathrm{slow}},
\label{eq:weight_update}
\end{equation}
where $\eta$ is the learning rate. 
The slow trace acts as a molecular memory buffer, accumulating STDP events over seconds and enabling a delayed dopaminergic signal to reinforce the relevant synapses retroactively.

\subsection{Biological stabilizers}

\subsubsection{Global Activity Regulator (GAR)}
Recurrent SNNs with external stimulation are prone to two pathological regimes: seizure-like runaway excitation and a quiescent coma. 
To maintain activity in a biologically realistic band (approximately 5--10Hz), OrganoidEnv employs a global activity regulator that monitors the average firing rate $\bar{f}(t)$ in windows of 50ms and applies interventions:
\begin{equation}
\text{Intervention}(t) = 
\begin{cases} 
\text{Poisson noise} & \text{if } \bar{f}(t) < f_{\text{coma}} \\
\text{Global reset} & \text{if } \bar{f}(t) > f_{\text{seizure}} \\
\text{none} & \text{otherwise}
\end{cases}
\label{eq:gar_intervention}
\end{equation} 

where $f_{\text{coma}} = 2$ Hz and $f_{\text{seizure}} = 50$ Hz, respectively.
This mechanism regularly rescues the silent networks and terminates the seizure cascades.

\subsubsection{Inhibitory homeostasis and structural plasticity}
To counteract the runaway potentiation of excitatory synapses, the baseline inhibitory weight $w_0^{\mathrm{inh}}(t)$ is scaled proportionally to the mean excitatory weight $\bar{w}^{\mathrm{exc}}(t)$:
\begin{equation}
w_0^{\mathrm{inh}}(t) = \alpha\, \bar{w}^{\mathrm{exc}}(t),
\end{equation}
with a scaling factor $\alpha > 1$ \cite{HomeostaticSTDP2010,Turrigiano2008SelfTuning}. 

In addition, structural plasticity is applied every 10 episodes: synapses whose weights fall below a threshold of 0.1 are pruned from the connectivity matrix, thereby reducing noise and computational overhead. 
This pruning complements dual-trace plasticity by removing persistently weak connections.

\subsection{RL agent and sensory--motor interface}

The outer agent is a duelling double DQN with prioritised experience replay (PER) and Hindsight Experience Replay (HER) \cite{Wang2016dueling,Andrychowicz2017HER,VanHasselt2016}. 
The agent observes a continuous-valued state vector as follows:
\begin{itemize}
    \item robot position and goal position in the 2D arena,
    \item proximity sensor readings near obstacles,
    \item a binary context bit indicating the active goal.
\end{itemize}
The action space consists of $8$ discrete commands corresponding to combinations of the four motor quadrants (for example, Up, Up-Right, Right, etc.). 
Each action is translated into a current injection pattern that targets the appropriate motor clusters.

At each environmental step, the agent selects an action, the SNN is stimulated accordingly, and the spiking dynamics are integrated for a fixed period. 
The resulting motor activity displaces a cursor in the 2D environment, and the environment returns a scalar reward and termination flag. 
The RL hyperparameters are summarised in Table~\ref{table:hyperparameters}.

\subsection{Environment and curriculum}

The environment was a $1.0 \times 1.0$ continuous 2D arena with a maximum of 80 steps per episode. 
The reward function combines a sparse goal reward $R_{\mathrm{goal}}$, penalties for collisions with obstacles, and a small step penalty to encourage efficient trajectory generation. 
Curriculum learning is used to gradually increase the task difficulty.

\begin{itemize}
    \item \textbf{Stage 1 (Basic navigation, episodes 0--99):} One goal, no obstacles, relatively large goal radius; $R_{\mathrm{goal}} = 0.15$.
    \item \textbf{Stage 2 (Obstacles, episodes 100--199):} Static circular and box obstacles added; $R_{\mathrm{goal}} = 0.10$.
    \item \textbf{Stage 3 (Multi-goal, episodes 200--349):} Two or more potential goals, selected by a context bit; $R_{\mathrm{goal}} = 0.10$.
    \item \textbf{Stage 4 (Full task, episodes 350--499):} Obstacles + multi-goal context with smaller goal radius; $R_{\mathrm{goal}} = 0.05$.
\end{itemize}

HER with $k=4$ is applied to failed episodes by treating the final cursor position as an alternative goal, providing an additional learning signal in sparse-reward regimes \cite{Andrychowicz2017HER}.

\begin{table}[h]
\centering
\caption{Key Hyperparameters}
\label{table:hyperparameters}
\begin{tabular}{llc}
\hline
\textbf{Component} & \textbf{Parameter} & \textbf{Value} \\ \hline
\textit{Spiking Neural} & Integration timestep $\Delta t$ & 0.1 ms \\
\textit{Network (SNN)}  & ATP recovery time $\tau_{\text{rec}}$ & 100 ms \\
                        & Fast trace decay $\tau_{\text{fast}}$ & 100 ms \\
                        & Slow trace decay $\tau_{\text{slow}}$ & 2500 ms \\
                        & Neuron count & 500 \\
                        & SDM sparsity & 10\% \\
                        & Prune threshold & 0.1 \\
                        & Coma threshold $f_{\text{coma}}$ & 2 Hz \\
                        & Seizure threshold $f_{\text{seizure}}$ & 50 Hz \\ \hline
\textit{RL Agent}       & Discount factor $\gamma$ & 0.99 \\
\textit{(Dueling DDQN)} & Target update rate $\tau$ & 0.005 \\
                        & N-step returns & 3 \\
                        & HER $k$ & 4 \\
                        & Action space & 8 (discrete) \\
                        & Steps per episode & 80 \\ \hline
\end{tabular}
\end{table}

\section{Results}

\subsection{Overall performance}

Across the 500-episode curriculum, OrganoidEnv achieved an overall success rate of 70.8\% (354/500 episodes). 
Success is defined as reaching a goal within the step budget without colliding with obstacles. 
The stage-wise performance is as follows:
\begin{itemize}
    \item Stage 1 (Basic): 79.0\% success, peaking near 100\% between episodes 60--100.
    \item Stage 2 (Obstacles): 97.0\% success.
    \item Stage 3 (Multi-goal): 58.0\% success.
    \item Stage 4 (Full task): 60.7\% success.
\end{itemize}

The marked improvement between stages 1 and 2 shows that the SNN can robustly incorporate obstacle information via SDM expansion. 
The 60.7\% success rate in Stage 4 (full task) reflects the representational bottleneck introduced by context-dependent routing with a static connectivity structure.

\subsection{Ablation studies}

To quantify the importance of architectural components, several ablations were performed in an identical environment with RL hyperparameters. 
Table~\ref{table:ablation} summarizes the results.

\begin{table}[h]
\centering
\caption{Ablation Study Results}
\label{table:ablation}
\begin{tabular}{l c p{4cm}}
\hline
\textbf{Configuration} & \textbf{Success (\%)} & \textbf{Effect} \\ \hline
Full OrganoidEnv & 70.8 & Baseline performance \\
No dual-trace STDP & 24.5 & Credit assignment failure \\
No motor mapping & 12.0 & Degraded control signals \\
No GAR (stabilizer) & 5.2 & Seizure/Coma states \\
No SDM expansion & $\approx$ 0 & Spatial aliasing \\ \hline
\end{tabular}
\end{table}

    \begin{figure*}
        \centering
        \includegraphics[width=0.75\linewidth]{fig2_ablation_study.png}
        \caption{Ablation study results at Month 4 curriculum stage. (A) Goal completion rates and (B) mean episode rewards for the full OrganoidEnv architecture versus configurations with SDM/morphology or homeostasis removed. Error bars represent standard deviation across N seeds. Note: goal completion counts shown are for the 100-episode Stage 2 window, not cumulative across all 500 episodes.}
        \end{figure*}

        \begin{figure*}
            \centering
            \includegraphics[width=1\linewidth]{fig3_month6_dashboard.png}
            \caption{Grand unification results across the four-stage curriculum. (A) Training reward curve with 20-episode moving average. (B) Per-stage success rates. (C) Goal completion by target context (Target A vs Target B). Dashed vertical lines mark curriculum stage transitions.}
            \end{figure*}

            

                \begin{figure}
                    \centering
                    \includegraphics[width=0.75\linewidth]{fig5_baseline_comparison.png}
                    \caption{Baseline comparison with statistical validation. OrganoidEnv (solid line, mean +/- std) against ANN, RL-only, and SimpleSNN baselines across training episodes.}
                    \end{figure}

                    \begin{figure}
                        \centering
                        \includegraphics[width=1\linewidth]{fig6_network_activity.png}
\caption{Neural activity during the final curriculum stage. (A) Spike raster plot showing sparse, distributed firing across 100 sampled neurons over 1000 ms. (B) Motor neuron activation patterns for the Up (blue) and Down (red) quadrants, demonstrating directional selectivity.}
                        \label{fig:placeholder2}
                    \end{figure}
Removing the dual-trace rule collapses the performance to 24.5\%, which is consistent with the difficulty of delayed credit assignment under pure STDP. 
Eliminating explicit motor quadrant mapping yielded only 12.0% success, indicating that direct sensory--motor topography is crucial for the RL to efficiently explore stimulation patterns. 
Without GAR or homeostasis, the network rapidly enters seizure or coma regimes, preventing the learning process.

\subsection{Baseline comparisons}

To contextualise performance, OrganoidEnv was compared with two non-biological baselines trained on the same environment and curriculum:

\begin{itemize}
    \item \textbf{ANN (MLP) policy:} A multilayer perceptron trained end-to-end with DQN.
    \item \textbf{RL-only (direct):} DQN controlling the cursor directly without an SNN.
\end{itemize}

\begin{table}[t]
\caption{Baseline Comparison}
\label{tab:baselines}
\centering
\begin{tabular}{l l l l}
\hline
\textbf{Method} & \textbf{Succ. (\%)} & \textbf{Ep. to 60\%} & \textbf{Final Reward} \\
\hline
ANN (MLP) & 98.2 $\pm$ 0.5 & 45 & 125.4 \\
RL-only (direct) & 92.4 $\pm$ 1.2 & 82 & 98.6 \\
OrganoidEnv & 70.8 $\pm$ 3.4 & 315 & 45.2 \\
\hline
\end{tabular}
\end{table}

As expected, ANN-based baselines converged faster and achieved higher asymptotic performance. 
These baselines serve as performance upper bounds for the navigation task, not as the primary point of comparison for a framework designed around biological constraints.

\subsection{Firing rates}

During the final curriculum stage, the network maintained an average firing rate of approximately 8.5Hz across neurones, remaining within the targeted biological band.

\section{Discussion}

The sparse firing regime (8.5Hz average) suggests potential energy benefits on neuromorphic hardware, though formal benchmarking is left for future work.

\subsection{Representational bottlenecks in context-dependent routing}

The 60.7% success ceiling in Stage 4 exposes a representational limitation of the fixed 500-neuron topology. 
Although the SDM layer reduces spatial aliasing, the network must route identical spatial inputs to distinct motor outputs, conditional on a context bit. 
With a static connectivity structure, different context-dependent trajectories interfere, leading to cross-context interference and performance saturation.

One potential extension is activity-dependent synaptogenesis, in which new synapses are formed between neuron pairs that exhibit persistent co-activity, analogous to BDNF-mediated axonal sprouting in biological tissues. 
Combined with pruning and homeostasis, this could allow the network to grow new pathways for novel contexts without catastrophic loss of knowledge.

\subsection{Position within neuromorphic and organoid intelligence}

OrganoidEnv is at the intersection of neuromorphic computing and organoid intelligence. 
It preserves key biological constraints, such as spikes, metabolic fatigue, local plasticity, and homeostatic stabilisers, while leveraging deep RL as an outer optimisation loop. 
This division of labour mirrors proposals in the OI literature, where symbolic AI or RL agents interface with organoids as adaptive substrates \cite{Smirnova2023OI,Hartung2023OIworkshop,Xiong2024BridgingAI}.

\subsection{Limitations}

Several limitations should be emphasised.

\begin{itemize}
    \item \textbf{Scale:} Real brain organoids contain on the order of $10^5$ to $10^6$ neurons with 3D spatial organization and diffusion-mediated signaling. The 500-neuron, 2D OrganoidEnv is thus a highly abstracted analogue; scaling studies and 3D extensions are planned for future work.
    \item \textbf{Dimensionality:} The network is embedded in a 2D task; real organoids are 3D structures with complex spatial constraints and diffusion gradients.
    \item \textbf{Synaptic transmission:} Synapses are modeled as instantaneous current injections, omitting stochastic release and receptor kinetics.
    \item \textbf{Glial cells and metabolism:} Glial regulation and detailed metabolic pathways are abstracted into a single ATP-like variable, potentially overestimating the efficacy of GAR relative to biological homeostasis.
    \item \textbf{Physics and hardware:} No explicit neuromorphic hardware mapping is implemented; energy estimates remain approximate.
\end{itemize}

These simplifications were deliberately performed to isolate the computational principles of RL-driven stimulation and local plasticity before introducing additional biological complexity.

\section{Conclusion}

This study presented OrganoidEnv, a stabilised RL environment in which a deep RL agent trains a biologically constrained SNN via current stimulation, while learning within the SNN is mediated by local dopamine-modulated dual-trace plasticity and homeostatic mechanisms. 
The system achieved robust performance on a continuous 2D navigation task with obstacles and context-dependent goals, and ablation studies confirmed the necessity of GAR, SDM expansion, dual-trace plasticity, and direct motor mapping. 
OrganoidEnv provides a controllable in silico platform for exploring neuromorphic and organoid-inspired computation and can serve as a bridge between theoretical SNN learning rules and future experiments with biological organoids.

\begin{thebibliography}{99}

\bibitem{Izhikevich2003simple}
E.~M. Izhikevich, ``Simple model of spiking neurons”, \emph{IEEE Trans. Neural Netw.}, vol.~14, no.~6, pp. 1569--1572, 2003.

\bibitem{NeuromorphicReview2019}
G.~Indiveri and S.-C. Liu, ``Neuromorphic computing: From materials to systems architecture,'' \emph{Nat. Rev. Mater.}, 2015.

\bibitem{NeuromorphicSNNMemristor2019}
S.~B. Furber, ``Neuromorphic spiking neural networks and their memristor implementations”, \emph{Front. Neurosci.}, 2019.

\bibitem{RecentFrontiersSNN2022}
W.~Pfeiffer and T.~P. Vogels, ``Recent advances and new frontiers in spiking neural networks”, \emph{Front. Neurosci.}, 2022.

\bibitem{DirectTrainingSNN2024}
S.~Guo et al., ``Direct training high-performance deep spiking neural networks”, \emph{Nat. Commun.}, 2024.

\bibitem{Florian2007rstdp}
R.~V. Florian, ``Reinforcement learning through modulation of spike-timing-dependent synaptic plasticity”, \emph{Neural Comput.}, vol.~19, no.~6, pp. 1468--1502, 2007.

\bibitem{FarriesFairhall2007}
M.~A. Farries and A.~L. Fairhall, ``Reinforcement learning with modulated STDP”, in \emph{Advances in Neural Information Processing Systems}, 2007.

\bibitem{HomeostaticSTDP2010}
A.~J. Watt and N.~S. Desai, ``Homeostatic plasticity and STDP: keeping a neuron’s cool in a fluctuating world,'' \emph{J. Neurosci.}, 2010.

\bibitem{Turrigiano2008SelfTuning}
G.~G. Turrigiano, ``The self-tuning neuron: synaptic scaling of excitatory synapses,'' \emph{Cell}, 2008.

\bibitem{eLifeHomeostatic2025}
H.~Lu et al., ``The interplay between homeostatic synaptic scaling and structural plasticity maintains firing rate homeostasis”, \emph{eLife}, 2025.

\bibitem{AcetylcholineNav2018}
R.~Frémaux et al., ``Acetylcholine-modulated plasticity in reward-driven navigation”, \emph{Sci. Rep.}, 2018.

\bibitem{Bellec2018l2l}
G.~Bellec et al., ``Long short-term memory and learning-to-learn in networks of spiking neurons,'' in \emph{Advances in Neural Information Processing Systems}, 2018.

\bibitem{Mnih2015human}
V.~Mnih et al., ``Human-level control through deep reinforcement learning”, \emph{Nature}, vol. 518, no. 7540, pp. 529--533, 2015.

\bibitem{VanHasselt2016}
H.~van Hasselt, A.~Guez, and D.~Silver, ``Deep reinforcement learning with double Q-learning”, in \emph{AAAI}, 2016.

\bibitem{Wang2016dueling}
Z.~Wang et al., ``Dueling network architectures for deep reinforcement learning”, in \emph{ICML}, 2016.

\bibitem{Andrychowicz2017HER}
M.~Andrychowicz et al., ``Hindsight experience replay”, in \emph{Advances in Neural Information Processing Systems}, 2017.

\bibitem{Smirnova2023OI}
L.~Smirnova et al., ``Organoid intelligence: the new frontier in biocomputing and intelligence-in-a-dish,'' \emph{Front. Sci.}, 2023.

\bibitem{Hartung2023OIworkshop}
T.~Hartung et al., ``First Organoid Intelligence (OI) workshop to form an OI community”, 2023.

\bibitem{Xiong2024BridgingAI}
J.~Xiong et al., ``Bridging artificial intelligence for biological computing and organoid intelligence”, IntechOpen, 2024.

\bibitem{SNNNavigation2023}
A.~Author et al., ``Navigation and the efficiency of spatial coding using reward-driven STDP in spiking neural networks”, \emph{bioRxiv}, 2023.

\bibitem{Davies2024Loihi}
M.~Davies et al., ``Advancing neuromorphic computing with Loihi 2,'' \emph{Nat. Electron.}, 2024.

\bibitem{Zheng2024SNNRL}
H.~Zheng et al., ``Deep reinforcement learning with spiking neural networks,'' \emph{IEEE Trans. Neural Netw. Learn. Syst.}, 2024.

\end{thebibliography}

\end{document}